\subsection{The limited  dependent variable models}
In many practical applications, it is common to encounter
phenomena with discrete or censored dependent variables.
Our \pkg{BACE} package includes several models that can
handle this issue, including the most popular ones such
as logit and probit models, as well as their ordered
specifications and multinomial logistic regression.
Additionally, the package provides support for Poisson
regression, Tobit model, and interval regression%
\footnote{The \pkg{BACE} package also includes linear
	regression models. However, since the dependent variable
	is not bounded in these models, linear regression is
	not the main focus of our interest.}.
In all of these cases, the dependent variable is either
dichotomous (i.e., taking the values zero or one),
categorical, a count, or censored.
In the remainder of this subsection, we present the standard
and popular models for limited dependent variable included
in our \pkg{BACE} package, which are also discussed in
\cite{koop2007bayesian, hosmer:logistic, Winkelmann_2008, verbeek2018guide}
and \cite{cottrell_luchcetti_gretl_manual}.

First, let us analyze the logit and probit models. The logit
and probit models are two types of regression models commonly
used to predict a binary dependent variable. Logit models are
based on the logistic distribution, while probit models are
based on the normal distribution. Let us consider a model of the form:
\begin{equation}
	z_i= \sum_{j=1}^{k_r}x_{ij}\beta_{j} + \epsilon_i, 
\end{equation}
and	
\begin{equation*}	
	y_i= \left \{\begin{array}{ll}
					1 & \text{if } z_i >0 \\
					0 & \text{if } z_i \leq 0,
				\end{array}\right.
\end{equation*}
where $z_i$, $i=1, \dots, N$ is an unobserved variable, while
$y_i$ is binary, observed data. $x_{ij}$ and $\beta_j$ are
explanatory variables and regression coefficients in model
$M_r$ for $j=1,\dots, k_r$.
If $\epsilon_i \overset{\text{iid}}{\sim} N(0, 1)$, then we
have a probit model. Otherwise, we obtain a logit regression.

In the next step, let us specify logit and probit models, where
we consider the choice between $L$ alternatives, numbered from 1 to $L$:
\begin{equation}
		z_i= \sum_{j=1}^{k_r}x_{ij}\beta_{j} + \epsilon_i,
%		z_i= x_{ri}\beta_{ri} + \epsilon_i, 
\end{equation}
and
\begin{equation*}
		y_i= \left \{ 
			\begin{array}{lc}
					 1 & \text{if } \alpha_0 < z_i \leq  \alpha_1 \\
					 2 & \text{if } \alpha_1 < z_i \leq  \alpha_2 \\
					 \vdots & \vdots \\
					L & \text{if } \alpha_{L-1} < z_i \leq  \alpha_{L}, \\
			\end{array}\right.
\end{equation*}
where unknown $\alpha_i$s are cutpoints with
$\alpha_0 = -\infty, \alpha_1 = 0$ and $\alpha_L = \infty$.
Assuming that $\epsilon_i$ is iid and standard normal, we obtain
the ordered probit model, whereas the logistic distribution
gives the ordered logit model.

The latest considered model of this kind is multinomial logistic
regression. It is an extension of logistic regression modeling
for cases where the dependent variable is not binary but can
take the values $0,1, \ldots, p$ and the categories are nominal,
meaning they do not have a natural order.

The model can be written in the following form:
\begin{align}
	\Prob\left(y_i= k\mid x_i\right) &= \frac{\exp \left( x_i\beta_k\right)} {1+\sum_{j=1}^{p}\exp{(x_i\beta_j)}}, k > 0,\\
	\Prob\left(y_i= 0\mid x_i\right) &= \frac{1} {1+\sum_{j=1}^{p}\exp{(x_i\beta_j)}}.
\end{align}

Poisson regression is a model used to predict count data
occurrences. That is, a dependent variable $y$ takes values
$k=0,1, 2\ldots$ from the Poisson distribution. Under this
assumption, the probability of the outcome variable $y$
conditional upon regressors is given by:
\begin{equation}
	\Prob\left(y_i= k\mid x_i\right)  = \frac{\exp \left\{-\lambda_{i} \right\} \lambda_{i}^k}{k!},
\end{equation}
where $\lambda_i= \exp(\sum_{j=1}^{k_r}x_{ij}\beta_{j})$ in model $M_r$.

The Tobit model is a type of regression model used to analyze
censored data, which refers to data with a lower or upper bound
on the dependent variable. Consider the following Tobit model
(type I) with two thresholds \citep{cottrell_luchcetti_gretl_manual}:
\begin{equation}
	z_i= \sum_{j=1}^{k_r}x_{ij}\beta_{j} + \epsilon_i,
	\label{eq: unobs_dep}
\end{equation}
and
\begin{equation}
	y_i= \left \{ 
	\begin{array}{ll}
		a & \text{for }  z_i \le a\\
		z_i &\text{for }  a < z_i < b\\
		b & \text{for }  z_i \ge b,\\
	\end{array}   \right.
\label{eq: tobit}
\end{equation}
where $\epsilon \sim N(0, \sigma^2)$. Usually the most
common specification is  $a=0$ and $b= \infty$.That is,
at zero, observations are censored from below, which means
that we do not observe any $z_i$ values that are negative \citep{verbeek2018guide}.

Another model that is available in the \pkg{BACE} package is
the interval regression. It is similar to the standard linear
regression, but it allows for the possibility of censoring
some data points due to measurement limitations or other
factors \citep{Stewart_1983, Bettin_Lucchetti_2012, cottrell_luchcetti_gretl_manual}.
In this case, we consider a data generating process with an
unobserved dependent variable $z_i$ (Equation \ref{eq: unobs_dep})
and an observed interval that contains this variable:
\begin{equation}
	m_i \leq z_i \leq  M_i.
\end{equation}
The interval can also be left-unbounded, i.e., $m_i = -\infty$
or right-unbounded, when $M_i = \infty$. In the case of
$m_i = M_i$, all observations are point observations and the
model reduces to a standard regression model, which is also
available in the \pkg{BACE} package. Please note that due to
the complexity of the censoring problem, the interval
regression model is only available in script form.

All of the mentioned models are available in the \pkg{BACE}
package. Through the use of model averaging analysis, it is
possible to estimate parameters, identify the key determinants,
and make predictions.

% BACE 3.2: interpretation and model averaging of ORs and AMEs.
\subsection{Model-averaged odds ratios and marginal effects}
\label{sec:or_ame_theory}
For binary and ordered response models, the regression coefficients
alone do not describe changes in outcome probabilities. The
\pkg{BACE} package therefore provides model-averaged odds ratios
(ORs) for binary and ordered logit models and average marginal
effects (AMEs) for binary and ordered logit and probit models.
Both quantities are first calculated within each model $M_r$ and
then averaged over the model space in the same way as the
parameter estimates in Section~\ref{MA_basics}. As there, the sums
over $2^K$ models are in practice approximated using the models
visited by the \MCcube algorithm.

We use the notation of the models described above. In model $M_r$
the latent variable is $z_i=x_i\beta+\epsilon_i=\sum_{j}x_{ij}\beta_j+\epsilon_i$,
where $x_{ij}$ is the $j$-th of $K$ candidate regressors and
$\beta_j=0$ whenever $x_j$ is not included in $M_r$. Let $F$ and $f$
denote the cumulative distribution function and the density of
$\epsilon_i$, i.e.\ the logistic ones for logit and the standard normal
ones for probit models. For a binary dependent variable
\begin{equation}
	\Prob(y_i=1\mid x_i,M_r)=F(x_i\beta)=1-\Prob(y_i=0\mid x_i,M_r),
	\label{eq:bace_prob_binary}
\end{equation}
and for an ordered one with $L$ categories and cutpoints
$\alpha_0=-\infty<\alpha_1<\dots<\alpha_L=\infty$
\begin{equation}
	\Prob(y_i=l\mid x_i,M_r)=F(\alpha_l-x_i\beta)-F(\alpha_{l-1}-x_i\beta),
	\quad l=1,\dots,L.
	\label{eq:bace_prob_ordered}
\end{equation}
Here $l$ indexes the observed outcome categories in increasing order;
their labels need not form a consecutive sequence.

\paragraph{Odds ratios for logit models.}
In the binary logit model the odds $\Prob(y_i=1\mid x_i,M_r)/\Prob(y_i=0\mid x_i,M_r)$
equal $\exp(x_i\beta)$, so a one-unit increase in $x_{ij}$ multiplies
the odds by $\mathrm{OR}_j=\exp(\beta_j)$. In the ordered logit model
(\ref{eq:bace_prob_ordered}), $\Prob(y_i\le l\mid x_i,M_r)=F(\alpha_l-x_i\beta)$,
so $\exp(\beta_j)$ is the common factor by which the odds
$\Prob(y_i>l\mid x_i,M_r)/\Prob(y_i\le l\mid x_i,M_r)$ change, the same
for every $l$. If $x_j\notin M_r$, then $\beta_j=0$ and $\mathrm{OR}_j=1$.

To take within-model coefficient uncertainty into account, we use the normal
approximation $\beta_j\mid y,M_r\ \dot{\sim}\ N\bigl(\E(\beta_j\mid y,M_r),\VAR(\beta_j\mid y,M_r)\bigr)$,
where $\E(\beta_j\mid y,M_r)$ and $\VAR(\beta_j\mid y,M_r)$ are the
MLE estimate and its estimated variance from model $M_r$. Then
$\mathrm{OR}_j$ is approximately lognormal and, for $x_j\in M_r$,
%\begin{align}
%	\E(\mathrm{OR}_j\mid y,M_r)&=\exp\!\Bigl(\E(\beta_j\mid y,M_r)+\tfrac12 \VAR(\beta_j\mid y,M_r)\Bigr),
%	\label{eq:bace_or_moment1}\\
%	\VAR(\mathrm{OR}_j\mid y,M_r)&=\Bigl[\exp\bigl(\VAR(\beta_j\mid y,M_r)\bigr)-1\Bigr]
%	\exp\!\Bigl(2\E(\beta_j\mid y,M_r)+\VAR(\beta_j\mid y,M_r)\Bigr),
%	\label{eq:bace_or_moment2}
%\end{align}
\begin{align}
	\E(\mathrm{OR}_j\mid y,M_r)&=\exp\!\Bigl(\E(\beta_j\mid y,M_r)+\tfrac12 \VAR(\beta_j\mid y,M_r)\Bigr),
	\label{eq:bace_or_moment1}\\
	\VAR(\mathrm{OR}_j\mid y,M_r)&=\Bigl[\exp\bigl(\VAR(\beta_j\mid y,M_r)\bigr)-1\Bigr]\nonumber\\
	&\quad\times\exp\!\Bigl(2\E(\beta_j\mid y,M_r)+\VAR(\beta_j\mid y,M_r)\Bigr),
	\label{eq:bace_or_moment2}
\end{align}
whereas for $x_j\notin M_r$ we have $\E(\mathrm{OR}_j\mid y,M_r)=1$ and
$\VAR(\mathrm{OR}_j\mid y,M_r)=0$. The model-averaged OR and its variance
are obtained analogously to $\E(\theta\mid y)$ and $\VAR(\theta\mid y)$:
\begin{align}
	\E(\mathrm{OR}_j\mid y)&\approx\sum_{r=1}^{2^K}\Prob(M_r\mid y)\,\E(\mathrm{OR}_j\mid y,M_r),\\
	\VAR(\mathrm{OR}_j\mid y)&\approx\sum_{r=1}^{2^K}\Prob(M_r\mid y)\,\VAR(\mathrm{OR}_j\mid y,M_r)\nonumber\\
	&\quad+\sum_{r=1}^{2^K}\Prob(M_r\mid y)\Bigl(\E(\mathrm{OR}_j\mid y,M_r)-\E(\mathrm{OR}_j\mid y)\Bigr)^2,
	\label{eq:bace_or_bma}
\end{align}
and the reported standard deviation is $\sqrt{\VAR(\mathrm{OR}_j\mid y)}$.
Thus, $\E(\mathrm{OR}_j\mid y)$ is \emph{not} simply the exponential of
the model-averaged coefficient $\E(\beta_j\mid y)$. In particular, an OR of one
corresponds to no change in odds, and the posterior mass of models
excluding a regressor contributes an OR of one. Exponentiating a
probit coefficient does not produce an odds ratio, so ORs are not
reported for probit models. These regressor-level ORs should not be
confused with the posterior \emph{model} odds ratios in
Equation~\ref{eq:posterior_odds_ratio}.

\paragraph{Average marginal effects.}
For a continuous regressor $x_j$, the average marginal effect on
the probability of category $l$ in model $M_r$ is
\begin{equation}
	\mathrm{AME}_{jl}(M_r)
	=\frac1{N}\sum_{i=1}^{N}
	\frac{\partial\Prob(y_i=l\mid x_i,M_r)}{\partial x_{ij}}.
	\label{eq:bace_ame_continuous}
\end{equation}
From (\ref{eq:bace_prob_binary}), in a binary logit or probit
\begin{equation}
	\frac{\partial\Prob(y_i=1\mid x_i,M_r)}{\partial x_{ij}}=\beta_j f(x_i\beta),
	\label{eq:bace_ame_binary}
\end{equation}
with the negative of this effect for $y_i=0$, and from
(\ref{eq:bace_prob_ordered}), in an ordered logit or probit
\begin{equation}
	\frac{\partial\Prob(y_i=l\mid x_i,M_r)}{\partial x_{ij}}
	=\beta_j\Bigl[f(\alpha_{l-1}-x_i\beta)-f(\alpha_l-x_i\beta)\Bigr],
	\label{eq:bace_ame_ordered}
\end{equation}
where the density at an infinite cutpoint ($\alpha_0$ or $\alpha_L$)
is zero.

For regressors taking only the values zero and one, the software
instead computes the average discrete change
\begin{equation}
	\mathrm{AME}_{jl}^{(0\to1)}(M_r)=\frac1{N}\sum_{i=1}^{N}
	\Bigl[\Prob(y_i=l\mid x_{ij}=1,x_{i,-j},M_r)
	-\Prob(y_i=l\mid x_{ij}=0,x_{i,-j},M_r)\Bigr],
	\label{eq:bace_ame_binary_regressor}
\end{equation}
where $x_{i,-j}$ denotes the remaining regressors, held at their observed
values. For a regressor absent from $M_r$, $\beta_j=0$ and all its marginal
effects are zero. Finally, the reported effects are averaged across models:
\begin{equation}
	\E(\mathrm{AME}_{jl}\mid y)\approx\sum_{r=1}^{2^K}\Prob(M_r\mid y)\,\mathrm{AME}_{jl}(M_r),
	\label{eq:bace_ame_bma}
\end{equation}
with $\mathrm{AME}_{jl}^{(0\to1)}(M_r)$ in place of $\mathrm{AME}_{jl}(M_r)$
for 0/1 regressors. The model-specific effects are evaluated at the
estimated coefficients $\E(\beta_j\mid y,M_r)$ and cutpoints; unlike
the OR variance (\ref{eq:bace_or_bma}), the procedure does not propagate
the within-model sampling variance and does not supply a standard deviation
of the AMEs. In every model and in the model average, the marginal effects
over the exhaustive outcome categories sum to zero (up to numerical
rounding), i.e.\ $\sum_{l}\E(\mathrm{AME}_{jl}\mid y)=0$. An effect of
$0.01$ corresponds to one percentage point.

Automatic recognition of 0/1 regressors requires both values to occur in
the estimation sample and is currently used only when no lagged dependent
variables have been introduced. Otherwise the program reports average
derivatives. For several dummy variables encoding one categorical regressor,
the 0-to-1 comparison is performed on each dummy separately; users should
interpret such changes with care, because a counterfactual holding the other
dummies fixed can represent an impossible combination of categories.
These are descriptive model-based effects, not causal effects without
additional identifying assumptions.

\subsection{Jointness measures}\label{sec:jointness_measures}
The main implementations of  model averaging are concerned
with the selection of variables when model uncertainty is
present. Another relevant issue which arises in this framework is to
identify whether different sets of two variables $x_i$ and $x_j$ are
substitutes, complements or neither over the model space. For that
reason, \cite{LeySteel2007} and \cite{DoppelhoferWeeks2009} define
ex-post jointness measures of dependence between different sets of
explanatory variables. Following \cite{Hofmarcher2018},
we can  use a general contingency Table \ref{tab:GeneralContingencyTable}
and easily calculate a number of jointness measures considered
in the literature. In \pkg{BACE} we have:

\begin{table}[h!tb]
	\centering
	\caption{General contingency table for two binary variables.}
	\begin{tabular}{l|ccc}\hline
		&\textit{B}	&	$\overline{B}$	&Total	\\\hline
		\textit{A}		&	a		&	b			& $a+b$		\\
		$\overline{A}$	&	c		&	d			& $c+d$		\\
		Total			&	$a+c$	&	$b+d$		& $a+b+c+d=N$\\\hline
	\end{tabular}
	\label{tab:GeneralContingencyTable}
\end{table}

\begin{itemize}\label{sec:jointess_numbers}
	\item The jointness measure $J_{\mathit{DW1}}$, first
		proposed in the \cite{DoppelhoferWeeks2005}:
		\begin{align}
			\begin{split}
				J_{\mathit{DW1}}(A,B) &  = \ln\left(\frac{\Prob(AB\mid y)}
				{\Prob(A\mid y) \Prob(B\mid y)}\right)\\
				& = \ln\left(\frac{aN}{(a+b)(a+c)}\right).
			\end{split}
			\label{DW_jointness1}
		\end{align}
	\item An alternative variant $J_{\mathit{DW2}}$ also
		proposed by \cite{DoppelhoferWeeks2009}:
		\begin{align}
			\begin{split}
				J_{\mathit{DW2}}(A,B) & =\ln\left(\frac {\Prob(AB \mid y)\Prob(\overline{A}\overline{B} \mid y)}
				{\Prob(A\overline{B}\mid y)\Prob(\overline{A}B\mid y)}\right) \\
				& = \ln\left(\frac{ad}{bc}\right).
			\end{split}
			\label{DW_jointness2}
		\end{align}
	\item The jointness statistics presented in \cite{LeySteel2007}:
		\begin{align}
			\begin{split}
				J_{\mathit{LS1}}(A,B) & = \frac{\Prob(AB\mid y)}{\Prob(A\mid y)+\Prob(B\mid y)-\Prob(AB\mid y)} \\
				& = \frac{a}{a+b+c}.
			\end{split}
			\label{LS_jointness1}
		\end{align}
	\item An additional variant of an earlier measure:
		\begin{align}
			\begin{split}
				J_{\mathit{LS2}}(A,B) & = \frac{\Prob(AB\mid y)}{\Prob(A\overline{B}\mid y)+\Prob(\overline{A}B\mid y)} \\
				& =  \frac{a}{b+c}.
			\end{split}
			\label{LS_jointness2}
		\end{align}
	\item The measure  proposed in \cite{Strachan2009}:
		\begin{align}
			\begin{split}
				J_{\mathit{St}}(A,B) & = \Prob(A\mid y)\Prob(B\mid y)
				\ln\left(\frac{\Prob(AB\mid y)}{\Prob(A\overline{B}\mid y)
					\Prob(\overline{A}B\mid y)}\right) \\
				& = \frac{a+b}{N}\frac{a+c}{N}\ln\left(\frac{aN}{bc}\right).
			\end{split}
			\label{St_jointness}
		\end{align}
	\item Yule's Q measure:
		\begin{align}
			\begin{split}
				J_{\mathit{YQ}}(A,B) & = \frac{\Prob(AB)\Prob(\overline{AB})-\Prob(A\overline{B})\Prob(\overline{A}B)}
				{\Prob(AB)\Prob(\overline{AB})+\Prob(A\overline{B})\Prob(\overline{A}B)} \\
				& = \frac{ad - bc}{ad + bc}.
			\end{split}
			\label{Yq_jointness}
		\end{align}
	\item The last  measure and also the default one,
		modified Yule's Q:
		\begin{equation}		
			J_{\mathit{YQ}} = \frac
			{\left(a+\frac{1}{2}\right)\left(d+\frac{1}{2}\right) -
				\left(b+\frac{1}{2}\right)\left(c+\frac{1}{2}\right)} 
			{\left(a+\frac{1}{2}\right)\left(d+\frac{1}{2}\right) +
				\left(b+\frac{1}{2}\right)\left(c+\frac{1}{2}\right) - \frac{1}{2}}.
			\label{YqM_jointness}
		\end{equation}
\end{itemize}

Information on the properties of these listed measures
and their interpretations can be found in the publication
of \cite{Hofmarcher2018}.
